% GENERATED FILE. Edit canonical lessons, then run npm run generate:books.
% canonical-source-hash: fnv1a64:c6630365e3572c6f
% canonical-lessons: TA-C307-ceytikkurippu, TA-C307-ceytit-totarpalar, TA-C307-pattirikai, TA-C307-tunukku, TA-C307-arivi

\chapter{Press Releases and Titbits}
\label{ch:ta-press-releases-and-titbits}
\hlchaptermodality{307}

\begin{chapteropening}
\textbf{By the end of this chapter:} \emph{I can say how news reaches the papers.}

Complete the last lesson of chapter 307: I can say how news reaches the papers.
\end{chapteropening}

\section[\texorpdfstring{\ta{செய்திக்குறிப்பு} (ceytikkuṟippu)}{செய்திக்குறிப்பு (ceytikkuṟippu)}]{\ta{செய்திக்குறிப்பு} (ceytikkuṟippu) — a press release}

\label{lesson:TA-C307-ceytikkurippu}



\begin{quote}
Before the new one: say the Tamil for an edition, then the Tamil for to bring out.
\end{quote}

\subsection*{You'll want to know: \ta{செய்திக்குறிப்பு}}
\textbf{\ta{செய்திக்குறிப்பு}} — \emph{ceytikkuṟippu} — \textquotedblleft{}a press release\textquotedblright{}.

\begin{quote}
\ta{காவல்துறை} \ta{ஒரு} \ta{செய்திக்குறிப்பை} \ta{வெளியிட்டது}.
\end{quote}

\begin{quote}
\emph{kāvaltuṟai oru ceytikkuṟippai veḷiyiṭṭatu.}
\end{quote}

\begin{quote}
The police issued a press release.
\end{quote}

\ta{செய்தி}, news + \ta{குறிப்பு}, a note.

\begin{cousinweb}
\textbf{\ta{செய்திக்குறிப்பு}} (\emph{ceytikkuṟippu}) is a Tamil compound built from \ta{செய்தி}, \textquotedblleft{}news\textquotedblright{} (DEDR 1957), and \ta{குறிப்பு}, \textquotedblleft{}a note\textquotedblright{}, from \ta{குறி}, \textquotedblleft{}to aim, mark\textquotedblright{} (DEDR 1847).
\end{cousinweb}

\begin{grammarlens}[title={What you've built}]
One more word for telling what happened.
\end{grammarlens}

\subsection*{Your turn}
\begin{itemize}
\raggedright
  \item \emph{Say it:} \emph{ceytikkuṟippu}
  \item \emph{Say it:} \emph{ceytikkuṟippu}, once more
  \item \emph{Say it:} \emph{kāvaltuṟai oru ceytikkuṟippai veḷiyiṭṭatu.} — The police issued a press release.
  \item \emph{Recall:} say the Tamil for an edition, then the Tamil for to bring out, then say \emph{ceytikkuṟippu} again
\end{itemize}

\begin{tcolorbox}[breakable,colback=teal!4,colframe=teal!35!black,title={Before you move on}]
What does \ta{செய்திக்குறிப்பு} mean? (\textquotedblleft{}A press release\textquotedblright{}.) Say it once more.
\end{tcolorbox}
\section[\texorpdfstring{\ta{செய்தித்} \ta{தொடர்பாளர்} (ceytit toṭarpāḷar)}{செய்தித் தொடர்பாளர் (ceytit toṭarpāḷar)}]{\ta{செய்தித்} \ta{தொடர்பாளர்} (ceytit toṭarpāḷar) — a spokesperson}

\label{lesson:TA-C307-ceytit-totarpalar}



\begin{quote}
Before the new one: say the Tamil for to bring out, then the Tamil for a press release.
\end{quote}

\subsection*{You'll want to know: \ta{செய்தித்} \ta{தொடர்பாளர்}}
\textbf{\ta{செய்தித்} \ta{தொடர்பாளர்}} — \emph{ceytit toṭarpāḷar} — \textquotedblleft{}a spokesperson\textquotedblright{}.

\begin{quote}
\ta{அரசின்} \ta{செய்தித்} \ta{தொடர்பாளர்} \ta{கேள்விகளுக்குப்} \ta{பதில்} \ta{சொன்னார்}.
\end{quote}

\begin{quote}
\emph{araciṉ ceytit toṭarpāḷar kēḷvikaḷukkup patil coṉṉār.}
\end{quote}

\begin{quote}
The government spokesperson answered questions.
\end{quote}

\ta{தொடர்பு}, contact + -\ta{ஆளர்}: the person in contact with the press.

\begin{cousinweb}
\textbf{\ta{செய்தித்} \ta{தொடர்பாளர்}} (\emph{ceytit toṭarpāḷar}) is a Tamil phrase built from \ta{செய்தி}, \textquotedblleft{}news\textquotedblright{} (DEDR 1957), and \ta{தொடர்பாளர்}, \textquotedblleft{}one who keeps the link\textquotedblright{}, from \ta{தொடர்}, \textquotedblleft{}to follow, be linked\textquotedblright{} (DEDR 3480).
\end{cousinweb}

\begin{grammarlens}[title={What you've built}]
One more word for telling what happened.
\end{grammarlens}

\subsection*{Your turn}
\begin{itemize}
\raggedright
  \item \emph{Say it:} \emph{ceytit toṭarpāḷar}
  \item \emph{Say it:} \emph{ceytit toṭarpāḷar}, once more
  \item \emph{Say it:} \emph{araciṉ ceytit toṭarpāḷar kēḷvikaḷukkup patil coṉṉār.} — The government spokesperson answered questions.
  \item \emph{Recall:} say the Tamil for to bring out, then the Tamil for a press release, then say \emph{ceytit toṭarpāḷar} again
\end{itemize}

\begin{tcolorbox}[breakable,colback=teal!4,colframe=teal!35!black,title={Before you move on}]
What does \ta{செய்தித்} \ta{தொடர்பாளர்} mean? (\textquotedblleft{}A spokesperson\textquotedblright{}.) Say it once more.
\end{tcolorbox}
\section[\texorpdfstring{\ta{பத்திரிகை} (pattirikai)}{பத்திரிகை (pattirikai)}]{\ta{பத்திரிகை} (pattirikai) — a magazine; the press}

\label{lesson:TA-C307-pattirikai}



\begin{quote}
Before the new one: say the Tamil for a press release, then the Tamil for a spokesperson.
\end{quote}

\subsection*{You'll want to know: \ta{பத்திரிகை}}
\textbf{\ta{பத்திரிகை}} — \emph{pattirikai} — \textquotedblleft{}a magazine; the press\textquotedblright{}.

\begin{quote}
\ta{என்} \ta{அப்பா} \ta{தினமும்} \ta{ஒரு} \ta{பத்திரிகை} \ta{படிக்கிறார்}.
\end{quote}

\begin{quote}
\emph{eṉ appā tiṉamum oru pattirikai paṭikkiṟār.}
\end{quote}

\begin{quote}
My father reads a magazine every day.
\end{quote}

Also the press in general: \ta{பத்திரிகைச்} \ta{செய்தி}, a press report.

\begin{cousinweb}
\textbf{\ta{பத்திரிகை}} (\emph{pattirikai}) is borrowed from Sanskrit \emph{pattrikā}, \textquotedblleft{}a leaf, a letter\textquotedblright{}, from \emph{pattra}, \textquotedblleft{}leaf, leaf of a book\textquotedblright{} (CDIAL 7733).
\end{cousinweb}

\begin{grammarlens}[title={What you've built}]
One more word for telling what happened.
\end{grammarlens}

\subsection*{Your turn}
\begin{itemize}
\raggedright
  \item \emph{Say it:} \emph{pattirikai}
  \item \emph{Say it:} \emph{pattirikai}, once more
  \item \emph{Say it:} \emph{eṉ appā tiṉamum oru pattirikai paṭikkiṟār.} — My father reads a magazine every day.
  \item \emph{Recall:} say the Tamil for a press release, then the Tamil for a spokesperson, then say \emph{pattirikai} again
\end{itemize}

\begin{tcolorbox}[breakable,colback=teal!4,colframe=teal!35!black,title={Before you move on}]
What does \ta{பத்திரிகை} mean? (\textquotedblleft{}A magazine; the press\textquotedblright{}.) Say it once more.
\end{tcolorbox}
\section[\texorpdfstring{\ta{துணுக்கு} (tuṇukku)}{துணுக்கு (tuṇukku)}]{\ta{துணுக்கு} (tuṇukku) — a titbit; a snippet}

\label{lesson:TA-C307-tunukku}



\begin{quote}
Before the new one: say the Tamil for a spokesperson, then the Tamil for a magazine.
\end{quote}

\subsection*{You'll want to know: \ta{துணுக்கு}}
\textbf{\ta{துணுக்கு}} — \emph{tuṇukku} — \textquotedblleft{}a titbit; a snippet\textquotedblright{}.

\begin{quote}
\ta{நாளிதழில்} \ta{ஒரு} \ta{வேடிக்கையான} \ta{துணுக்கைப்} \ta{படித்தேன்}.
\end{quote}

\begin{quote}
\emph{nāḷitaḻil oru vēṭikkaiyāṉa tuṇukkaip paṭittēṉ.}
\end{quote}

\begin{quote}
I read a funny titbit in the newspaper.
\end{quote}

A short piece in a paper; at root a small bit of anything.

\begin{cousinweb}
\textbf{\ta{துணுக்கு}} (\emph{tuṇukku}) is native Dravidian (DEDR 3305): a small piece, from \ta{துணி}, \textquotedblleft{}to be cut, severed\textquotedblright{}.
\end{cousinweb}

\begin{grammarlens}[title={What you've built}]
One more word for telling what happened.
\end{grammarlens}

\subsection*{Your turn}
\begin{itemize}
\raggedright
  \item \emph{Say it:} \emph{tuṇukku}
  \item \emph{Say it:} \emph{tuṇukku}, once more
  \item \emph{Say it:} \emph{nāḷitaḻil oru vēṭikkaiyāṉa tuṇukkaip paṭittēṉ.} — I read a funny titbit in the newspaper.
  \item \emph{Recall:} say the Tamil for a spokesperson, then the Tamil for a magazine, then say \emph{tuṇukku} again
\end{itemize}

\begin{tcolorbox}[breakable,colback=teal!4,colframe=teal!35!black,title={Before you move on}]
What does \ta{துணுக்கு} mean? (\textquotedblleft{}A titbit; a snippet\textquotedblright{}.) Say it once more.
\end{tcolorbox}
\section[\texorpdfstring{\ta{அறிவி} (aṟivi)}{அறிவி (aṟivi)}]{\ta{அறிவி} (aṟivi) — to announce}

\label{lesson:TA-C307-arivi}



\begin{quote}
Before the new one: say the Tamil for a magazine, then the Tamil for a titbit.
\end{quote}

\subsection*{You'll want to know: \ta{அறிவி}}
\textbf{\ta{அறிவி}} — \emph{aṟivi} — \textquotedblleft{}to announce\textquotedblright{}.

\begin{quote}
\ta{அரசு} \ta{நாளை} \ta{விடுமுறை} \ta{என்று} \ta{அறிவித்தது}.
\end{quote}

\begin{quote}
\emph{aracu nāḷai viṭumuṟai eṉṟu aṟivittatu.}
\end{quote}

\begin{quote}
The government announced that tomorrow is a holiday.
\end{quote}

Past (I …): \ta{அறிவித்தேன்}, \emph{aṟivittēṉ}.

What is announced comes before \ta{என்று}; the noun is \ta{அறிவிப்பு}.

\begin{cousinweb}
\textbf{\ta{அறிவி}} (\emph{aṟivi}) is native Dravidian: DEDR 314 gives it beside \ta{அறி}, \textquotedblleft{}to know\textquotedblright{}, so to announce is to make known.
\end{cousinweb}

\begin{grammarlens}[title={What you've built}]
One more word for telling what happened.
\end{grammarlens}

\subsection*{Your turn}
\begin{itemize}
\raggedright
  \item \emph{Say it:} \emph{aṟivi}
  \item \emph{Say it:} \emph{aṟivi}, once more
  \item \emph{Say it:} \emph{aracu nāḷai viṭumuṟai eṉṟu aṟivittatu.} — The government announced that tomorrow is a holiday.
  \item \emph{Recall:} say the Tamil for a magazine, then the Tamil for a titbit, then say \emph{aṟivi} again
\end{itemize}

\begin{tcolorbox}[breakable,colback=teal!4,colframe=teal!35!black,title={Before you move on}]
What does \ta{அறிவி} mean? (\textquotedblleft{}To announce\textquotedblright{}.) Say it once more.
\end{tcolorbox}
